On introduit dans un erlenmeyer~:
\begin{itemize}
    \item $V_1=\qty{25.0}{\mL}$ d'acide propanoïque $\ce{C2H5CO2H_{(aq)}}$ de
          concentration~: ${C_1=\qty{1.0e-2}{\mol\per\L}}$~;
    \item \raggedright $V_2=\qty{25.0}{\mL}$ de propanoate de sodium
          ${\ce{Na^{+}_{(aq)} + C2H5CO^{-}_{2(aq)}}}$ de concentration~:
          ${C_2=\qty{2.5e-2}{\M}}$~;
    \item $V_3=\qty{30.0}{\mL}$ de solution de chlorure d'hydroxylammonium
          $\ce{NH3OH^{+}_{(aq)} + Cl^{-}_{(aq)}}$, de
          concentration~: $C_3=\qty{1.5e-2}{\mol\per\L}$~; 
    \item $V_4=\qty{20.0}{\mL}$ d'hydroxylamine $\ce{NH2OH_{(aq)}}$ de
          concentration~: $C_4=\qty{1.0e-2}{\mol\per\L}$.
\end{itemize}
\begin{enumerate}
    \item Écrire son équation en considérant que l'acide propanoïque est un
          réactif.
    \item Exprimer et calculer la constante d'équilibre $\mathrm{K}$ de cette
          réaction.
    \item Déterminer la valeur du quotient de réaction $Q_{r,i}$ dans l'état
          initial de ce système.
    \item Indiquer, en justifiant, dans quel sens évolue spontanément le système
          chimique à partir de son état initial.
\end{enumerate}

\begin{donnees}
    \begin{itemize}
        \item $\mathrm{pK_{A}}(\ce{C2H5CO2H}\ttslash{}\ce{C2H5CO^-_2})=
              \mathrm{pK_{A_1}}=4,8$~; \
              $\mathrm{pK_{A}}(\ce{NH3OH^{+}}\ttslash{}\ce{NH2OH})=
              \mathrm{pK_{A_2}}=6,1$.
    \end{itemize}
\end{donnees}
